\documentclass{beamer}

\usepackage[T1]{fontenc}
\usepackage[utf8]{inputenc}
\usepackage{lmodern}
\usepackage{amsmath,amssymb,amsfonts}
\usepackage{physics}
\usepackage{bm}
\usepackage{mathtools}
\usepackage{quantikz}

\title{Quantum Fourier Transform and Quantum Phase estimation algorithms}
\author{Morten Hjorth-Jensen}
\date{FYS5419/9419: Week of March 23-27, 2026}

\begin{document}

%------------------------------------------------
\frame{\titlepage}

%------------------------------------------------
\begin{frame}{Outline}
\begin{enumerate}
\item From the classical DFT to the QFT, reminder from last week
\item Circuit decomposition with Hadamards and controlled phases
\item Approximate QFT and complexity
\item Quantum Phase Estimation (QPE)
\item Connection to Hamiltonian simulation
\item Connection to Variational Quantum Eigensolvers (VQE)
\item Themes for the second project
\end{enumerate}
See also separate jupyter-notebooks at \url{https://github.com/CompPhysics/QuantumComputingMachineLearning/tree/gh-pages/doc/pub/week10/ipynb}
\end{frame}



%------------------------------------------------
\begin{frame}{Motivation}
The Quantum Fourier Transform (QFT) is one of the central algorithms in quantum computing. 

It appears in:
\begin{itemize}
\item Shor's factoring algorithm,
\item order finding,
\item period finding,
\item quantum phase estimation,
\item several quantum simulation algorithms (HHL and other).
\end{itemize}

The QFT is the quantum analogue of the discrete Fourier transform, but it acts on the amplitudes of a quantum state.
\end{frame}

%------------------------------------------------
\begin{frame}{Classical Discrete Fourier Transform, reminder from last week}
Given a vector \(x=(x_0,x_1,\dots,x_{N-1})\), the discrete Fourier transform is
\[
y_k = \frac{1}{\sqrt N}\sum_{j=0}^{N-1} x_j e^{2\pi i jk/N},
\qquad k=0,\dots,N-1.
\]

Equivalently, the DFT matrix is
\[
F_N = \frac{1}{\sqrt N}
\begin{pmatrix}
1 & 1 & 1 & \cdots & 1 \\
1 & \omega & \omega^2 & \cdots & \omega^{N-1} \\
1 & \omega^2 & \omega^4 & \cdots & \omega^{2(N-1)} \\
\vdots & \vdots & \vdots & \ddots & \vdots \\
1 & \omega^{N-1} & \omega^{2(N-1)} & \cdots & \omega^{(N-1)^2}
\end{pmatrix}
\]
with
\[
\omega = e^{2\pi i/N}.
\]
\end{frame}

%------------------------------------------------
\begin{frame}{Definition of the QFT}
For \(N=2^n\), the QFT is the unitary transformation
\[
\operatorname{QFT}_N \ket{x}
=
\frac{1}{\sqrt N}
\sum_{k=0}^{N-1}
e^{2\pi i xk/N}\ket{k},
\qquad x=0,\dots,N-1.
\]

Thus the QFT maps a computational basis state \(\ket{x}\) to a coherent superposition of all basis states with phases determined by \(x\).

Because the map is linear, for a general state
\[
\ket{\psi}=\sum_{x=0}^{N-1} c_x \ket{x},
\]
we have
\[
\operatorname{QFT}_N \ket{\psi}
=
\sum_{x=0}^{N-1} c_x \operatorname{QFT}_N\ket{x}.
\]
\end{frame}

%------------------------------------------------
\begin{frame}{Binary Representation of the Input}
Let
\[
x = x_1 x_2 \dots x_n
\]
be the binary representation of the integer \(x\), meaning
\[
x = \sum_{j=1}^n x_j 2^{n-j},
\qquad x_j\in\{0,1\}.
\]

We also write binary fractions as
\[
0.x_j x_{j+1}\dots x_n
=
\frac{x_j}{2}+\frac{x_{j+1}}{2^2}+\cdots+\frac{x_n}{2^{n-j+1}}.
\]

This notation is crucial because the QFT phases factor naturally into such binary fractions.
\end{frame}

%------------------------------------------------
\begin{frame}{Step-by-Step Derivation of the Product Form I}
Start from
\[
\operatorname{QFT}_N \ket{x}
=
\frac{1}{\sqrt{2^n}}
\sum_{k=0}^{2^n-1}
e^{2\pi i xk/2^n}\ket{k}.
\]

Write the output label \(k\) in binary:
\[
k = k_1 k_2 \dots k_n,
\qquad
k = \sum_{m=1}^n k_m 2^{n-m}.
\]

Then
\[
\frac{xk}{2^n}
=
x \sum_{m=1}^n \frac{k_m}{2^m}.
\]

Hence
\[
e^{2\pi i xk/2^n}
=
\prod_{m=1}^n e^{2\pi i x k_m / 2^m}.
\]
\end{frame}

%------------------------------------------------
\begin{frame}{Step-by-Step Derivation of the Product Form II}
Because each \(k_m\in\{0,1\}\), the sum over all \(k\) factorizes:
\[
\operatorname{QFT}_N \ket{x}
=
\frac{1}{2^{n/2}}
\bigotimes_{m=1}^n
\left(
\ket{0} + e^{2\pi i x/2^m}\ket{1}
\right).
\]

Now only the fractional part of \(x/2^m\) matters in the phase, and that fractional part is exactly
\[
0.x_{n-m+1}x_{n-m+2}\dots x_n.
\]

Thus
\[
\operatorname{QFT}_N \ket{x_1 x_2 \dots x_n}
=
\frac{1}{2^{n/2}}
\left(\ket{0}+e^{2\pi i\,0.x_n}\ket{1}\right)
\otimes
\left(\ket{0}+e^{2\pi i\,0.x_{n-1}x_n}\ket{1}\right)
\otimes \cdots
\]
\[
\cdots \otimes
\left(\ket{0}+e^{2\pi i\,0.x_1x_2\dots x_n}\ket{1}\right).
\]

The output order is therefore reversed relative to the input qubit order.
\end{frame}

%------------------------------------------------
\begin{frame}{Final Product Form}
The standard product form is
\[
\operatorname{QFT}_N \ket{x_1x_2\dots x_n}
=
\frac{1}{2^{n/2}}
\left(\ket{0}+e^{2\pi i\,0.x_n}\ket{1}\right)
\otimes
\left(\ket{0}+e^{2\pi i\,0.x_{n-1}x_n}\ket{1}\right)
\otimes \cdots
\]
\[
\cdots \otimes
\left(\ket{0}+e^{2\pi i\,0.x_1x_2\dots x_n}\ket{1}\right).
\]

This factorization is the key reason the QFT can be implemented efficiently.

It shows that the QFT can be built from:
\begin{itemize}
\item Hadamard gates,
\item controlled phase rotations,
\item a final qubit reversal using SWAP gates.
\end{itemize}
\end{frame}

%------------------------------------------------
\begin{frame}{Basic Gates Used in the QFT}
The main one-qubit gate is the Hadamard
\[
H = \frac{1}{\sqrt 2}
\begin{pmatrix}
1 & 1 \\
1 & -1
\end{pmatrix},
\]
which creates equal superpositions.

The phase rotations are
\[
R_k =
\begin{pmatrix}
1 & 0 \\
0 & e^{2\pi i / 2^k}
\end{pmatrix}.
\]

Controlled-\(R_k\) gates apply \(R_k\) to a target qubit only if the control qubit is \(\ket{1}\).
\end{frame}

%------------------------------------------------
\begin{frame}{How the Circuit Arises}
The first output qubit in the product form contains the phase
\[
e^{2\pi i\,0.x_1x_2\dots x_n}.
\]

This is produced by:
\begin{itemize}
\item a Hadamard on the first qubit, generating \(\frac{1}{\sqrt2}(\ket0+\ket1)\),
\item then a sequence of controlled phase gates from the later qubits, adding the binary-fraction phase.
\end{itemize}

The same pattern repeats recursively for the remaining qubits.

Thus the QFT circuit is naturally built qubit by qubit.
\end{frame}

%------------------------------------------------
\begin{frame}{Two-Qubit QFT: Mathematical Form}
For \(n=2\), we have \(N=4\), and
\[
\operatorname{QFT}_4 \ket{x_1x_2}
=
\frac{1}{2}
\left(\ket0 + e^{2\pi i\,0.x_2}\ket1\right)
\otimes
\left(\ket0 + e^{2\pi i\,0.x_1x_2}\ket1\right).
\]

This requires:
\begin{itemize}
\item Hadamard on the first qubit,
\item controlled-\(R_2\),
\item Hadamard on the second qubit,
\item then a SWAP to correct the reversed output order.
\end{itemize}
\end{frame}

%------------------------------------------------
\begin{frame}{Two-Qubit QFT Circuit}
\centering
\begin{quantikz}
\lstick{$q_0$} & \gate{H} & \ctrl{1} & \qw & \swap{1} & \qw \\
\lstick{$q_1$} & \qw & \gate{R_2} & \gate{H} & \targX{}{} & \qw
\end{quantikz}


Here:
\begin{itemize}
\item \(H\) creates the superposition,
\item \(R_2=\mathrm{diag}(1,e^{i\pi/2})\),
\item the SWAP reverses the qubit order.
\end{itemize}
\end{frame}

%------------------------------------------------
\begin{frame}{Three-Qubit QFT: Mathematical Structure}
For \(n=3\),
\[
\operatorname{QFT}_8 \ket{x_1x_2x_3}
=
\frac{1}{\sqrt 8}
\left(\ket0+e^{2\pi i\,0.x_3}\ket1\right)
\otimes
\left(\ket0+e^{2\pi i\,0.x_2x_3}\ket1\right)
\otimes
\left(\ket0+e^{2\pi i\,0.x_1x_2x_3}\ket1\right).
\]

The first qubit therefore requires phase contributions corresponding to
\[
\frac{x_2}{2} + \frac{x_3}{4},
\]
and this is exactly what the controlled \(R_2\) and \(R_3\) gates generate.
\end{frame}

%------------------------------------------------
\begin{frame}{Three-Qubit QFT Circuit}
\centering
\begin{quantikz}
\lstick{$q_0$} & \gate{H} & \ctrl{1}    & \ctrl{2}    & \qw      & \qw        & \qw      & \swap{2}   & \qw \\
\lstick{$q_1$} & \qw      & \gate{R_2}  & \qw         & \gate{H} & \ctrl{1}   & \qw      & \qw        & \qw \\
\lstick{$q_2$} & \qw      & \qw         & \gate{R_3}  & \qw      & \gate{R_2} & \gate{H} & \targX{}{} & \qw
\end{quantikz}

\vspace{0.5em}

Conceptually:
\begin{enumerate}
\item \(H\) on qubit 0, then controlled \(R_2,R_3\)
\item \(H\) on qubit 1, then controlled \(R_2\)
\item \(H\) on qubit 2
\item final reversal of qubit order
\end{enumerate}
\end{frame}

%------------------------------------------------
\begin{frame}{General \(n\)-Qubit QFT Circuit}
For qubit \(j\), counting from the top:

\begin{enumerate}
\item Apply a Hadamard \(H\)
\item For each later qubit \(k>j\), apply a controlled \(R_{k-j+1}\)
\item Continue recursively
\item At the end, reverse qubit order with SWAP gates
\end{enumerate}

The general structure is:
\[
H \to \text{controlled phases} \to H \to \text{controlled phases} \to \cdots \to \text{SWAP network}.
\]
\end{frame}

%------------------------------------------------
\begin{frame}{General Pattern in Quantikz}
\centering
\begin{quantikz}
\lstick{$q_0$} & \gate{H} & \ctrl{1} & \ctrl{2} & \ctrl{3} & \qw \\
\lstick{$q_1$} & \qw & \gate{R_2} & \ctrl{1} & \ctrl{2} & \qw \\
\lstick{$q_2$} & \qw & \qw & \gate{R_2} & \ctrl{1} & \qw \\
\lstick{$q_3$} & \qw & \qw & \qw & \gate{R_2} & \gate{H}
\end{quantikz}

\vspace{0.5em}

This is a schematic pattern rather than a fully expanded full-\(n\) circuit. The key point is the nested structure of Hadamards and controlled phase gates.
\end{frame}

%------------------------------------------------
\begin{frame}{Worked 3-Qubit Phase Example I}
Take the computational basis input
\[
\ket{x}=\ket{101}.
\]

This means
\[
x_1=1,\qquad x_2=0,\qquad x_3=1,
\qquad x=5.
\]

The product form gives
\begin{align*}
\operatorname{QFT}_8 \ket{101}
= &
\frac{1}{\sqrt8}
\left(\ket0 + e^{2\pi i\,0.1}\ket1\right)
\otimes
\left(\ket0 + e^{2\pi i\,0.01}\ket1\right) \\
\otimes &
\left(\ket0 + e^{2\pi i\,0.101}\ket1\right).
\end{align*}
\end{frame}

%------------------------------------------------
\begin{frame}{Worked 3-Qubit Phase Example II}
Now evaluate the binary fractions:
\[
0.1 = \frac12,
\qquad
0.01 = \frac14,
\qquad
0.101 = \frac12 + \frac18 = \frac58.
\]

Hence
\[
e^{2\pi i\,0.1} = e^{i\pi} = -1,
\]
\[
e^{2\pi i\,0.01} = e^{i\pi/2} = i,
\]
\[
e^{2\pi i\,0.101} = e^{2\pi i(5/8)} = e^{5\pi i/4}.
\]

Therefore
\[
\operatorname{QFT}_8 \ket{101}
=
\frac{1}{\sqrt8}
(\ket0-\ket1)\otimes(\ket0+i\ket1)\otimes(\ket0+e^{5\pi i/4}\ket1).
\]
\end{frame}

%------------------------------------------------
\begin{frame}{Worked 3-Qubit Phase Example III}
This example shows exactly how the QFT encodes the binary digits of the input into relative phases.

This phase encoding is what makes the QFT so useful in:
\begin{itemize}
\item phase estimation,
\item period finding,
\item extracting eigenphases of unitaries.
\end{itemize}

In other words, the QFT converts arithmetic structure into measurable interference structure.
\end{frame}

%------------------------------------------------
\begin{frame}{Inverse QFT}
The inverse QFT is obtained by reversing the circuit and taking Hermitian conjugates:
\begin{itemize}
\item replace each \(R_k\) by \(R_k^\dagger\),
\item reverse the order of the gates,
\item keep the same SWAP network.
\end{itemize}

Thus the inverse QFT is also efficient.

In practice, many algorithms use the inverse QFT rather than the forward QFT.
\end{frame}

%------------------------------------------------
\begin{frame}{Complexity of the QFT}
For \(n\) qubits:
\begin{itemize}
\item \(n\) Hadamard gates,
\item \(\frac{n(n-1)}{2}\) controlled phase gates,
\item about \(\lfloor n/2 \rfloor\) SWAP gates.
\end{itemize}

Therefore the total gate count is
\[
O(n^2).
\]

This is exponentially more efficient than applying the classical DFT matrix directly to a \(2^n\)-dimensional vector. The latter requires $O(N^2)$ operations, with $N=2^n$. With three qubits we have already $9$ versus $64$ operations. With 100 qubits we have $10^4$ compared with $2^{200}\sim 1.6\times 10^{60}$.  
\end{frame}

%------------------------------------------------
\begin{frame}{Approximate QFT}
Very small-angle phase gates contribute little.

Hence we may truncate the circuit by dropping controlled \(R_k\) gates for large \(k\).

This gives the approximate QFT:
\begin{itemize}
\item fewer gates,
\item shallower circuit,
\item often still sufficient for practical algorithms.
\end{itemize}

With suitable truncation, the gate complexity can be reduced to approximately
\[
O(n \log n).
\]
\end{frame}


%------------------------------------------------
\begin{frame}{Implementation}
A typical implementation of the QFT proceeds as follows:

\begin{enumerate}
\item Choose the number of qubits \(n\)
\item Loop over qubits \(j=0,\dots,n-1\)
\item Apply \(H\) on qubit \(j\)
\item For each later qubit \(k>j\), apply a controlled phase \(R_{k-j+1}\)
\item At the end, reverse the qubit order with SWAP gates
\end{enumerate}

This is exactly the constructive logic derived from the product form.
See the jupyter-notebook at \url{https://github.com/CompPhysics/QuantumComputingMachineLearning/tree/gh-pages/doc/pub/week10} for implementaions.
\end{frame}

%------------------------------------------------
\begin{frame}{Pseudocode for QFT Construction}
\[
\begin{array}{l}
\texttt{for } j=0 \texttt{ to } n-1: \\
\qquad \texttt{apply H on qubit } j \\
\qquad \texttt{for } k=j+1 \texttt{ to } n-1: \\
\qquad\qquad \texttt{apply controlled-}R_{k-j+1}\texttt{ from } k \texttt{ to } j \\
\texttt{for } j=0 \texttt{ to } \lfloor n/2\rfloor-1: \\
\qquad \texttt{swap qubit } j \texttt{ with qubit } n-1-j
\end{array}
\]

Depending on software conventions, the control and target order may appear reversed relative to the mathematical derivation, but the logical structure is the same.
\end{frame}

%------------------------------------------------
\begin{frame}{Simulation Remarks}
When simulating the QFT numerically, it is useful to compare:
\begin{itemize}
\item the circuit output state,
\item the exact matrix formula
\[
\operatorname{QFT}_N\ket{x}
=
\frac{1}{\sqrt N}\sum_{k=0}^{N-1} e^{2\pi i xk/N}\ket{k},
\]
\item and the product-form prediction.
\end{itemize}

This provides a clean consistency check:
\begin{itemize}
\item amplitudes should have equal magnitude \(1/\sqrt N\),
\item phases should match the Fourier kernel,
\item output qubit order must be checked carefully.
\end{itemize}
\end{frame}

%------------------------------------------------
\begin{frame}{Common Source of Confusion: Bit Reversal}
A frequent issue is the ordering of bits.

The product form naturally produces the qubits in reversed order:
\[
(x_1x_2\dots x_n)\mapsto
(\text{phase involving }x_n)\otimes \cdots \otimes (\text{phase involving }x_1\dots x_n).
\]

Thus:
\begin{itemize}
\item either include final SWAP gates,
\item or accept the reversed output convention.
\end{itemize}

Many software libraries omit the final SWAPs and simply document the reversed ordering.
\end{frame}

%------------------------------------------------
\begin{frame}{Why the QFT Matters Physically}
The QFT is not just an abstract transform.

It provides a mechanism for:
\begin{itemize}
\item converting phase information into computational-basis information,
\item exploiting interference to extract periodicity,
\item connecting time evolution phases to measurable outputs.
\end{itemize}

In this sense, the QFT is a bridge between
\[
\text{dynamical phase accumulation}
\qquad\longrightarrow\qquad
\text{algorithmic readout}.
\]
\end{frame}

%------------------------------------------------
\begin{frame}{Summary}
The Quantum Fourier Transform:
\begin{itemize}
\item is the quantum analogue of the discrete Fourier transform,
\item acts as
\[
\ket{x}\mapsto \frac{1}{\sqrt N}\sum_k e^{2\pi i xk/N}\ket{k},
\]
\item factorizes into one-qubit phase states,
\item is implemented using \(H\), controlled \(R_k\), and SWAP gates,
\item has gate complexity \(O(n^2)\),
\item is central to phase estimation and many other algorithms.
\end{itemize}
\end{frame}


%================================================
\begin{frame}{Quantum Phase Estimation: Problem Statement}
Suppose \(U\) is a unitary operator and \(\ket{\psi}\) is one of its eigenstates:
\[
U\ket{\psi}=e^{2\pi i \phi}\ket{\psi},
\qquad 0\le \phi <1.
\]

The goal of Quantum Phase Estimation (QPE) is to estimate the phase \(\phi\).

Equivalently, QPE extracts the eigenvalue of \(U\) in phase form.
\end{frame}

%================================================
\begin{frame}{Why QPE Matters}
QPE is central because many physical problems can be mapped to eigenvalue estimation.

Examples:
\begin{itemize}
\item energy estimation in quantum chemistry,
\item eigenvalue problems in many-body physics,
\item Hamiltonian simulation,
\item Shor's algorithm and period finding.
\end{itemize}

If we can encode a Hamiltonian into a unitary, then QPE can extract its eigenvalues.
\end{frame}

%================================================
\begin{frame}{Registers Used in QPE}
QPE uses two registers:

\begin{itemize}
\item a control register of \(n\) qubits, initialized in \(\ket{0}^{\otimes n}\),
\item a system register, initialized in the eigenstate \(\ket{\psi}\).
\end{itemize}

So the initial state is
\[
\ket{0}^{\otimes n}\ket{\psi}.
\]

The control register stores the estimated phase bits.
\end{frame}

%================================================
\begin{frame}{Step 1: Create a Uniform Superposition}
Apply Hadamard gates to all control qubits:
\[
\ket{0}^{\otimes n}
\mapsto
\frac{1}{2^{n/2}}
\sum_{k=0}^{2^n-1}\ket{k}.
\]

So the total state becomes
\[
\frac{1}{2^{n/2}}
\sum_{k=0}^{2^n-1}
\ket{k}\ket{\psi}.
\]

This prepares the control register as a coherent superposition of all binary numbers \(k\).
\end{frame}

%================================================
\begin{frame}{Step 2: Controlled Powers of \(U\)}
Now apply controlled powers of \(U\):
\[
U^{2^0}, U^{2^1}, U^{2^2}, \dots, U^{2^{n-1}}.
\]

Each control qubit determines whether the corresponding power is applied.

Because \(\ket{\psi}\) is an eigenstate,
\[
U^{2^j}\ket{\psi} = e^{2\pi i 2^j \phi}\ket{\psi}.
\]

Hence each controlled unitary writes phase information into the control register.
\end{frame}

%================================================
\begin{frame}{Step 2: State After Controlled Powers}
If the control register basis state is \(\ket{k}\), then the controlled operations apply \(U^k\), giving
\[
U^k\ket{\psi} = e^{2\pi i k\phi}\ket{\psi}.
\]

Therefore the combined state becomes
\[
\frac{1}{2^{n/2}}
\sum_{k=0}^{2^n-1}
e^{2\pi i k\phi}\ket{k}\ket{\psi}.
\]

The system register factors out, and the control register now contains the phase information.
\end{frame}

%================================================
\begin{frame}{Key Observation}
The control register is in the state
\[
\frac{1}{2^{n/2}}
\sum_{k=0}^{2^n-1} e^{2\pi i k\phi}\ket{k},
\]
which is exactly a Fourier-type state.

This is why the inverse QFT is the correct next step:
it converts this phase-encoded superposition into a computational basis approximation of \(\phi\).
\end{frame}

%================================================
\begin{frame}{Step 3: Apply the Inverse QFT}
Applying \(\mathrm{QFT}^{-1}\) to the control register gives
\[
\mathrm{QFT}^{-1}
\left(
\frac{1}{2^{n/2}}
\sum_{k=0}^{2^n-1} e^{2\pi i k\phi}\ket{k}
\right).
\]

If
\[
\phi = \frac{m}{2^n}
\]
exactly for some integer \(m\), then this becomes
\[
\ket{m}.
\]

So the control register stores the exact \(n\)-bit binary expansion of \(\phi\).
\end{frame}

%================================================
\begin{frame}{Exact QPE Case}
Suppose
\[
\phi = 0.\phi_1\phi_2\dots\phi_n
=
\frac{m}{2^n}.
\]

Then after inverse QFT, the control register becomes
\[
\ket{\phi_1\phi_2\dots\phi_n}.
\]

Thus measuring the control register returns the exact binary digits of \(\phi\).

This is the cleanest case of QPE.
\end{frame}

%================================================
\begin{frame}{Step-by-Step Derivation with Binary Fractions}
Let
\[
\phi = 0.\phi_1\phi_2\dots\phi_n\dots
\]

Then one can show that the control-register state after controlled powers factorizes as
\[
\frac{1}{2^{n/2}}
\left(\ket0 + e^{2\pi i\, 2^{n-1}\phi}\ket1\right)
\otimes
\left(\ket0 + e^{2\pi i\, 2^{n-2}\phi}\ket1\right)
\otimes \cdots
\]
\[
\cdots \otimes
\left(\ket0 + e^{2\pi i\, \phi}\ket1\right).
\]

Because
\[
2^{n-1}\phi = \phi_1.\phi_2\phi_3\dots,
\]
only the fractional parts matter in the phases, so each qubit encodes one successive binary suffix.
\end{frame}

%================================================
\begin{frame}{How the Bits Emerge}
Using the binary expansion of \(\phi\), the control-register state becomes
\[
\frac{1}{2^{n/2}}
\left(\ket0 + e^{2\pi i\,0.\phi_n}\ket1\right)
\otimes
\left(\ket0 + e^{2\pi i\,0.\phi_{n-1}\phi_n}\ket1\right)
\otimes \cdots
\]
\[
\cdots \otimes
\left(\ket0 + e^{2\pi i\,0.\phi_1\phi_2\dots\phi_n}\ket1\right),
\]
which has exactly the structure of a QFT state.

Therefore applying \(\mathrm{QFT}^{-1}\) yields the bit string
\[
\phi_1\phi_2\dots\phi_n.
\]
\end{frame}

%================================================
\begin{frame}{QPE Circuit}
\centering
\begin{quantikz}
\lstick{$\ket{0}$}    & \gate{H} & \ctrl{3} & \qw      & \qw      & \gate[3]{\mathrm{QFT}^\dagger} & \meter{} \\
\lstick{$\ket{0}$}    & \gate{H} & \qw      & \ctrl{2} & \qw      & \ghost{\mathrm{QFT}^\dagger}   & \meter{} \\
\lstick{$\ket{0}$}    & \gate{H} & \qw      & \qw      & \ctrl{1} & \ghost{\mathrm{QFT}^\dagger}   & \meter{} \\
\lstick{$\ket{\psi}$} & \qw      & \gate{U^4} & \gate{U^2} & \gate{U} & \qw                          & \qw
\end{quantikz}

\vspace{0.5em}

This example shows a 3-qubit control register and one system register.
\end{frame}

%================================================
\begin{frame}{Worked Example: \(\phi=5/8\)}
Suppose
\[
\phi = \frac{5}{8} = 0.101.
\]

Then the ideal QPE output on a 3-qubit control register is
\[
\ket{101}.
\]

Indeed,
\[
e^{2\pi i \phi} = e^{2\pi i (5/8)}.
\]

After the controlled powers and inverse QFT, the control register collapses to the bit string \(101\) with probability 1.
\end{frame}

%================================================
\begin{frame}{When \(\phi\) Is Not Exactly Representable}
If \(\phi\) is not exactly of the form \(m/2^n\), then QPE does not return one exact bit string.

Instead, the output is a probability distribution peaked near the best \(n\)-bit approximation:
\[
\phi \approx \frac{m}{2^n}.
\]

The measurement outcome gives the nearest binary approximation with high probability.
\end{frame}

%================================================
\begin{frame}{Precision of QPE}
With \(n\) control qubits, QPE resolves phases to approximately
\[
\Delta \phi \sim \frac{1}{2^n}.
\]

Thus the precision improves exponentially with the number of control qubits.

This exponential precision is one of the main reasons QPE is so powerful.
\end{frame}

%================================================
\begin{frame}{From Phases to Energies}
Suppose \(H\) is a Hamiltonian with eigenstate \(\ket{E}\):
\[
H\ket{E} = E\ket{E}.
\]

Define the time-evolution operator
\[
U(t)=e^{-iHt}.
\]

Then
\[
U(t)\ket{E} = e^{-iEt}\ket{E}.
\]

Comparing with
\[
U\ket{\psi}=e^{2\pi i \phi}\ket{\psi},
\]
we identify
\[
\phi = -\frac{Et}{2\pi} \mod 1.
\]

Thus QPE can estimate energies if we can implement \(e^{-iHt}\).
\end{frame}

%================================================
\begin{frame}{Connection to Hamiltonian Simulation}
QPE requires controlled applications of
\[
U^{2^j} = e^{-iH t 2^j}.
\]

So QPE depends directly on the ability to perform Hamiltonian simulation.

In practice this may be done using:
\begin{itemize}
\item Trotter-Suzuki decompositions,
\item linear-combination-of-unitaries methods,
\item qubitization,
\item problem-specific simulation techniques.
\end{itemize}

Thus:
\[
\text{Hamiltonian simulation} + \text{QPE}
\quad \Rightarrow \quad
\text{energy estimation}.
\]
\end{frame}

%================================================
\begin{frame}{Physical Interpretation}
QPE measures the dynamical phase accumulated by an energy eigenstate under time evolution.

The algorithm turns the phase
\[
e^{-iEt}
\]
into a bit string encoding the energy.

So one can view QPE as a highly precise interferometric energy-measurement protocol.
\end{frame}

%================================================
\begin{frame}{Why QPE Is Important in Quantum Chemistry}
In quantum chemistry and many-body physics, the main goal is often to compute eigenenergies of a Hamiltonian.

If we can prepare an approximate eigenstate with sufficient overlap with the true ground state, then QPE can project onto the exact eigenvalue with high precision.

This makes QPE a natural route to:
\begin{itemize}
\item ground-state energies,
\item excited-state energies,
\item spectral properties.
\end{itemize}
\end{frame}

%================================================
\begin{frame}{Connection to VQE}
The Variational Quantum Eigensolver (VQE) and QPE aim at similar physics goals, but in different ways.

\textbf{VQE:}
\begin{itemize}
\item variational,
\item hybrid quantum-classical,
\item optimized by minimizing
\[
E(\bm{\theta}) = \bra{\psi(\bm{\theta})} H \ket{\psi(\bm{\theta})},
\]
\item suitable for noisy near-term devices.
\end{itemize}

\textbf{QPE:}
\begin{itemize}
\item not variational,
\item requires coherent controlled time evolution,
\item provides direct eigenvalue estimation,
\item suited to fault-tolerant quantum computers.
\end{itemize}
\end{frame}

%================================================
\begin{frame}{How VQE and QPE Complement Each Other}
A common long-term strategy is:
\begin{enumerate}
\item Use a variational method such as VQE to prepare a good approximate eigenstate
\[
\ket{\psi_{\mathrm{trial}}}
\]
with large overlap with the true eigenstate.

\item Use that state as the input to QPE.

\item QPE then refines the energy estimate to high precision.
\end{enumerate}

So VQE can be viewed as a state-preparation front end for QPE.
\end{frame}

%================================================
\begin{frame}{Overlap Requirement}
If the initial state is
\[
\ket{\psi_{\mathrm{trial}}} = \sum_j c_j \ket{E_j},
\]
then QPE outputs \(E_j\) with probability
\[
|c_j|^2.
\]

Therefore the success probability depends on the overlap with the desired eigenstate.

This is exactly why improved ans\"atze in VQE are also valuable for QPE-based workflows.
\end{frame}

%================================================
\begin{frame}{QPE vs VQE in Practice}
\textbf{VQE}
\begin{itemize}
\item shallow circuits,
\item many repeated measurements,
\item classical optimization loop,
\item noise tolerant,
\item limited by optimization difficulty and shot cost.
\end{itemize}

\textbf{QPE}
\begin{itemize}
\item deep coherent circuits,
\item requires controlled \(e^{-iHt}\),
\item very precise,
\item ideal for fault-tolerant hardware,
\item often viewed as the long-term gold standard for energy estimation.
\end{itemize}
\end{frame}

%================================================
\begin{frame}{Semiclassical Interpretation}
One may interpret QPE as repeatedly extracting successive binary digits of a phase.

The inverse QFT acts globally, but conceptually it decodes the phase one bit at a time.

This bitwise structure is one reason QPE is so elegant mathematically.
\end{frame}

%================================================
\begin{frame}{Implementation}
A implementation of the QPE typically proceeds as follows:
\begin{enumerate}
\item choose the number of control qubits,
\item initialize the system register in an eigenstate or approximate eigenstate,
\item apply Hadamards to the control register,
\item apply controlled \(U^{2^j}\),
\item apply inverse QFT,
\item measure the control register,
\item interpret the measurement as an estimate of \(\phi\).
\end{enumerate}

If the notebook also implements the QFT explicitly, it will usually reuse the same Hadamard + controlled-phase + swap pattern.
\end{frame}

%================================================
\begin{frame}{Summary of the Full Picture}
\[
\text{Hamiltonian } H
\quad \longrightarrow \quad
U(t)=e^{-iHt}
\quad \longrightarrow \quad
\text{QPE}
\quad \longrightarrow \quad
E
\]

The QFT is the phase-decoding engine inside QPE.

VQE and QPE are complementary:
\begin{itemize}
\item VQE is a near-term variational energy estimator,
\item QPE is a long-term high-precision eigenvalue estimator,
\item VQE can provide input states for QPE.
\end{itemize}
\end{frame}

%================================================
\begin{frame}{Summary}
In these notes we have seen that:
\begin{itemize}
\item the QFT maps
\[
\ket{x}\mapsto \frac{1}{\sqrt N}\sum_k e^{2\pi i xk/N}\ket{k},
\]
\item QPE uses controlled powers of a unitary and the inverse QFT to extract eigenphases,
\item for Hamiltonian evolution
\[
U(t)=e^{-iHt},
\]
QPE becomes an energy-estimation algorithm,
\item VQE and QPE address related eigenvalue problems with different hardware requirements.
\end{itemize}
\end{frame}

%================================================
\begin{frame}{Suggestions for the second project}
\begin{itemize}
\item Continuation of project 1 with  more realistic systems and using adaptive VQE on actual (real) quantum computers
\item Implementation and studies of the Quantum Approximate Optimization Algorithm (QAOA): Simulation of linear algebra systems on quantum computers, solving for example differential equations
\item Implementation of the HHL algorithm (needs QPE)
\item Applications and implementations of quantum machine learning algorithms
\item Implement quantum neural networks for PINNs
\item Detailed quantum chemistry applications.
\end{itemize}
\end{frame}



\begin{frame}{Suggestions for the second project}
\begin{itemize}
\item Studies of entanglement and physical realization of quantum gates and circuits
\item Implementing Shor's algorithm or other algorithms
\item Parallelizing the Variational Quantum Eigensolver with multi-GPU scaling, see \url{https://arxiv.org/abs/2601.09951}
\item Implementing and comparing the quantum phase estimation algorithm with VQE, see for example recent QPE paper at \url{"https://arxiv.org/abs/2601.05788}
\item Studies of hyperparameters in VQE, see for example \url{https://arxiv.org/abs/2601.16679}
\item Implementing unitary coupled cluster theory on a quantum computer, see \url{https://arxiv.org/abs/2501.15893}
\end{itemize}
\end{frame}

\begin{frame}{What is the QAOA algorithm?}
\begin{block}{QAOA stands for the Quantum Approximate Optimization Algorithm}
Each part of the name reflects something essential about the method:
\begin{itemize}
\item Quantum: it runs on a quantum computer, using qubits and quantum gates.
\item Approximate: it typically finds near-optimal solutions rather than exact ones.
\item Optimization: it is designed to solve optimization problems (like MaxCut, scheduling, portfolio selection).
\item Algorithm: it is a well-defined hybrid quantum–classical procedure.
\end{itemize}
\end{block}

QAOA is a variational quantum algorithm that uses a parameterized quantum circuit to search for good solutions to hard optimization problems.
\end{frame}

\begin{frame}{What is the HHL algorithm?}

The HHL algorithm—named after Aram Harrow, Avinatan Hassidim, and Seth
Lloyd—is one of the central quantum algorithms for solving linear
systems of equations.

The core idea is to solve a linear system
\[
A \mathbf{x} = \mathbf{b}.
\]
\end{frame}
\end{document}

